(1) 由于哈密顿量仅含非对角项，其本征能量 \(E_\pm(k)=\pm|t_1+t_2e^{ika}|=\pm\sqrt{t_1^2+t_2^2+2t_1t_2\cos ka}\)。区边 \(k=\pi/a\) 处 \(|E|=|t_2-t_1|\)，能隙宽度是 \(2|t_2-t_1|\)，仅当 \(t_1=t_2\) 闭合。对角元为零意味着手征对称 \(\sigma_zH(k)\sigma_z=-H(k)\)，从而能谱正负对称。

(2) 定义非对角复函数 \(q(k)=t_1+t_2e^{-ika}\)。随 \(k\) 绕一整个布里渊区，\(q(k)\) 在复平面上描出以 \(t_1\) 为圆心、\(t_2\) 为半径的圆。若 \(t_2>t_1\)，圆包含原点，绕数大小为1；若 \(t_2<t_1\)，原点在圆外，绕数为0。绕数正负还依赖 \(k\) 方向与子晶格约定，物理可观测的非平凡与平凡分类不依赖这个符号。

(3) 在半无限左端寻求 \(E=0\) 态，令所有 \(B_n\) 分量为零，则 \(A_n\) 的定态方程自动满足；对每个 \(B_n\) 的 Schrödinger 方程要求 \(t_1 A_n+t_2 A_{n+1}=0\)，故 \(A_{n+1}=-(t_1/t_2)A_n\)。其解 \(A_n=\mathcal N(-t_1/t_2)^{n-1}\)，可归一化当且仅当 \(t_2>t_1\)，归一化系数 \(\mathcal N=\sqrt{1-(t_1/t_2)^2}\)。

(4) 把振幅写成 \(|A_n|\propto e^{-(n-1)a/\xi}\)，得到 \(\xi=a/\ln(t_2/t_1)\)。当 \(t_2\to t_1^+\)，局域长度发散，与体能隙闭合一致。有限长度链可同时存在右边界态，左右尾部重叠约按 \((t_1/t_2)^N\) 缩放，导致两边缘态以指数小能量成对偏离零；只有半无限或满足额外子晶格数目不平衡的情形才保证所构造的精确零能。若加入破坏手征对称的在位势，零能钉扎一般也不再受保护。
